% Generated draft: requires mathematical and source review.
\documentclass[11pt]{article}
\usepackage[margin=25mm]{geometry}
\usepackage{fontspec}
\setmainfont{Latin Modern Roman}
\usepackage{amsmath,amssymb,mathtools,unicode-math}
\setmathfont{Latin Modern Math}
\usepackage{enumitem,needspace,xurl,hyperref}
\setlength{\parindent}{0pt}
\setlength{\parskip}{5pt}
\setlength{\emergencystretch}{3em}
\newcommand{\N}{\mathbb{N}}
\newcommand{\Z}{\mathbb{Z}}
\newcommand{\Q}{\mathbb{Q}}
\newcommand{\R}{\mathbb{R}}
\newcommand{\C}{\mathbb{C}}
\newcommand{\F}{\mathbb{F}}
\newcommand{\A}{\mathbb{A}}
\newcommand{\PP}{\mathbb{P}}
\newcommand{\E}{\mathbb{E}}
\newcommand{\eps}{\varepsilon}
\newcommand{\dd}{\,\mathrm{d}}
\begin{document}
\section*{UnsolvedMath ID: 30001238}
\textbf{Source problem code:} catalogue entry
\subsection{Definitions and mathematical statement}
For a conventional precise reading, let \(n\in\Z_{\geq 0}\), let \(X\) be a smooth irreducible complex algebraic variety of dimension \(n\), and let \(\mathfrak a\subsetneq\mathcal O_X\) be a nonzero coherent ideal sheaf. These conventions concerning varieties and ideals are explicit interpretive choices: the supplied statement does not specify them or whether dimension zero is included.

A prime divisor over \(X\) is a prime divisor \(E\) on a smooth variety \(Y\) equipped with a proper birational morphism \(f:Y\to X\). Write \(\operatorname{ord}_E(\mathfrak a)\) for the vanishing order of \(\mathfrak a\mathcal O_Y\) along \(E\), and define its log discrepancy by
\[
A_X(E)=1+\operatorname{ord}_E(K_Y-f^*K_X).
\]
The global log canonical threshold is
\[
\operatorname{lct}_X(\mathfrak a)
=\inf_{\substack{E\text{ a prime divisor over }X\\ \operatorname{ord}_E(\mathfrak a)>0}}
\frac{A_X(E)}{\operatorname{ord}_E(\mathfrak a)}.
\]
Under these conventions, set
\[
\mathcal T_n=
\left\{\operatorname{lct}_X(\mathfrak a):
\begin{array}{l}
X\text{ is a smooth irreducible complex algebraic variety},\\
\dim X=n,\quad 0\ne\mathfrak a\subsetneq\mathcal O_X
\text{ is a coherent ideal sheaf}
\end{array}\right\}.
\]
A subset \(S\subseteq\R\) satisfies the ascending chain condition if every sequence \((c_i)_{i\geq1}\) in \(S\) with \(c_i\leq c_{i+1}\) is eventually constant; equivalently, \(S\) has no infinite strictly increasing sequence.

\textbf{Question recorded under permanent ID 30001238.} Is it true that, for every \(n\in\Z_{\geq0}\), the set \(\mathcal T_n\) satisfies the ascending chain condition? Equivalently, for each fixed \(n\), is there no sequence of pairs \((X_i,\mathfrak a_i)_{i\geq1}\) as above such that
\[
\operatorname{lct}_{X_1}(\mathfrak a_1)
<\operatorname{lct}_{X_2}(\mathfrak a_2)
<\operatorname{lct}_{X_3}(\mathfrak a_3)<\cdots\,?
\]
Here the dimension is fixed throughout each sequence, while both the variety and the ideal may vary. With the stated conventions, \(\mathcal T_0\) is empty.

\textbf{Unresolved source qualification.} The threshold question appears consistently in the supplied original, clean, and main statements. However, the supplied background and literature assessment concern a generalization of a Serre-condition regularity bound from simplicial and monomial settings. They report that the cited work described that other question as unknown even for monomial ideals, and that no complete resolution was located. No regularity-bound formulation is supplied. This assessment cannot be identified with an assessment of the threshold question from the supplied evidence. The intended association between the statement, report citation, and literature assessment therefore requires review.
\subsection{Short English statement}
For each fixed dimension, the recorded question asks whether log canonical thresholds of ideals on smooth complex varieties can form an infinite strictly increasing sequence. The displayed formulation uses explicit conventional definitions; the supplied background discusses a different regularity problem, so the source association and recorded status require review.
\subsection{Sources}
\begin{itemize}
\item Supplied local UnsolvedMath registry record, permanent ID \(30001238\), title \emph{Ascending Chains of Log Canonical Thresholds}, exact source problem code \texttt{catalogue entry}, in the \emph{Oberwolfach Reports Open Problems} collection. The registry assigns proposed year 2009; no individual proposer is supplied.
\item Supplied report attribution: \emph{Kommutative Algebra} (2009), Oberwolfach Reports, DOI 10.4171/owr/2009/22, pp.~1157--1224. \url{https://doi.org/10.4171/owr/2009/22}. The supplied background describes these reports as documenting mathematical workshops and open-problem discussions. No precise page or passage for this question is supplied.
\item Supplied literature source, labeled ``Murai--Terai on H-vectors and Serre conditions'': \url{https://arxiv.org/abs/0912.1089}.
\item Supplied literature source, labeled ``published Murai--Terai paper'': \url{https://doi.org/10.4310/MRL.2009.v16.n6.a10}.
\item Literature assessment checked on \texttt{2026-08-21}: this is supplied dated curation, not independently verified. It assigns status ``open'' and discusses a Serre-condition regularity-bound generalization, reportedly unknown even for monomial ideals in the cited work, with no complete resolution located. The record itself qualifies this as dated literature triage derived from supplied Oberwolfach curation and requests independent verification of time-sensitive status and cited claims. Its relevance to the threshold question is unresolved.
\item The supplied registry labels the statement \texttt{corrected\_verified} and asserts that its mathematical content was checked against the official Oberwolfach report, with notation, context, and question form normalized without changing the source-level claim. This is a supplied verification claim, not verification performed for this draft. The supplied ``open'' label is likewise recorded attribution, not an independently established mathematical or literature conclusion.
\end{itemize}
\end{document}
